\documentclass[10pt,a4paper]{amsart}
\usepackage[margin=19mm]{geometry}
\usepackage{amsmath,amssymb,mathtools}
\usepackage[scr=rsfs,cal=euler]{mathalfa}
\usepackage{xfrac}
\usepackage[unicode,hidelinks,bookmarks=false]{hyperref}
\newcommand{\ie}{i.e.\ }
\newcommand{\cf}{cf.\ }
\newcommand{\resp}{respectively\ }
\newcommand{\Subsection}{\subsection}
\providecommand{\nobreakdash}{\nobreak\hbox{-}}
\setlength{\emergencystretch}{2em}
\multlinegap0pt
\usepackage{amssymb}
\usepackage{xcolor}
\usepackage{algorithm}\usepackage{algpseudocode}
\usepackage{mathtools}
\newtheorem{lemma}{Lemma}
\usepackage{cleveref}
\crefname{lemma}{Lemma}{Lemmas}
\newcommand{\R}{\mathbb{R}}
\DeclareMathOperator*{\argmax}{arg\,max}
\DeclareMathOperator*{\argmin}{arg\,min}
\newcommand{\ip}[2]{\langle #1, #2 \rangle}
\title{Exact Sequence Interpolation with Transformers}
\author{Albert Alcalde, Giovanni Fantuzzi, Enrique Zuazua}
\date{Source: arXiv:2502.02270v3 (CC BY 4.0)}
\begin{document}
\maketitle

\noindent\emph{Source and licence.} Exact Sequence Interpolation with Transformers. Version arXiv:2502.02270v3 (CC BY 4.0), distributed by arXiv under the Creative Commons Attribution 4.0 International licence, http://creativecommons.org/licenses/by/4.0/, as recorded in the arXiv licence metadata for this exact version and shown on the abstract page https://arxiv.org/abs/2502.02270v3. Full licence notice: \texttt{licenses/arxiv-2502.02270-CC-BY-4.0.txt}.\par\smallskip

\noindent\textit{Editorial scope.} Counted unit: the article's lemma lem:rank1-asymptotics (source0.tex lines 561--569). The article's complete original proof of it is delivered with it (source0.tex lines 570--625, one \texttt{proof} environment of 56 lines). No mathematics is written, repaired or re-derived: the statement and the proof are the article's own lines, in the article's own notation, and no retained line is substituted or re-flowed.  Retained uncounted as the source-local context the proof needs: lines 545--548.  Nothing was omitted from any retained span, so the package carries no omission record.  No retained line contains a citation, so the wrapper carries no bibliography and no bibliography entry.  No retained line contains a figure environment, an image-inclusion command or drawing code, so no image is delivered and the package declares no asset dependency.  Editorial formatting only, in addition: an AMS \texttt{amsart} preamble that restates the article's own declarations and macros used by the retained lines, namely the environments \texttt{lemma} on separate counters, together with the standard packages those lines need.  The retained environments print in this document as Lemma 1 in order of appearance, whereas the article prints the same items with the numbering of its own full document.  The complete native source of the article as distributed in the arXiv e-print for this exact version is bundled unchanged as \texttt{sources/172719/source0.tex}.
\begin{equation}\label{eq:hardmaxDynamics}
    x_i(k+1) = \rho x_i(k) + \frac{1}{|\mathcal{C}_i(X(k), A)|} V
        \sum_{\ell \in \mathcal{C}_i(X(k), A)} x_\ell(k),
\end{equation}

\begin{lemma}\label{lem:rank1-asymptotics}
Given an initial configuration $\{ x_i(0) \}_{i \in [n]}$, suppose $v\in \R^d$ satisfies $\ip{v}{x_i(0)} \neq 0$ for all $i \in [n]$, and assume there exist unique indices $m,M\in [n]$ such that $\ip{v}{x_M(0)} > 0$, $\ip{v}{x_m(0)} < 0$ and 
$$
    \ip{v}{x_M(0)} > \ip{v}{x_j(0)} \quad \forall j \neq M,
    \qquad 
    \ip{v}{x_m(0)} < \ip{v}{x_j(0)} \quad \forall j \neq m.
$$
Then, the dynamics governed by \eqref{eq:hardmaxDynamics} with $\rho = 1 - \gamma$, $V= \gamma I_d$, $\gamma\in (0,1]$, and $A = vv^\top$ converge to an equilibrium consisting of $n_1$ copies of $x_m(0)$ and $n_2$ copies of $x_M(0)$ where $n_1 + n_2 = n$. For $\gamma=1$, convergence happens in one step.
\end{lemma}

\begin{proof}
If $\gamma=1$, the result follows from a direct calculation using  \eqref{eq:hardmaxDynamics} with $\rho=0$ and $V= I_d$.
If $\gamma < 1$, we start by proving that the following statements hold for all $k\geq 0$:
\begin{align}
&\argmax_{j\in[n]} \langle v, x_j(k)\rangle = \{M\}, \label{eq:(a)} \\
&\argmin_{j\in[n]} \langle v, x_j(k)\rangle = \{m\}, \label{eq:(b)} \\
&x_M(k) = x_M(0) \quad \text{and} \quad x_m(k) = x_m(0), \label{eq:(c)} \\[1em]
&\operatorname{sign}(\langle v, x_i(k)\rangle) = 
\operatorname{sign}(\langle v, x_i(0)\rangle) 
\quad \text{for all } i\in[n]. \label{eq:(d)}
\end{align}
These properties hold at $k=0$ by the assumptions. We now assume they hold at time $k$ and prove them at time $k+1$. First, we observe that
$\langle A x_i(k), x_j(k)\rangle
= \langle v, x_i(k)\rangle \langle v, x_j(k)\rangle$,
so
\[
\mathcal{C}_i(X(k),A)
=
\begin{cases}
\{M\}, & \langle v, x_i(k)\rangle > 0,\\
\{m\}, & \langle v, x_i(k)\rangle < 0.
\end{cases}
\]
This is because maximizing $\langle v, x_i(k)\rangle\langle v, x_j(k)\rangle$ over $j$ amounts to maximizing
$\langle v, x_j(k)\rangle$ if $\langle v, x_i(k)\rangle>0$ and minimizing it if $\langle v, x_i(k)\rangle<0$, and by the induction hypothesis the unique maximizer and minimizer are $M$ and $m$.

Next, we compute $x_i(k+1)$ using the update rule \eqref{eq:hardmaxDynamics} with $\rho=1-\gamma$ and $V=\gamma I_d$, and take the inner product with $v$ to find that
\begin{equation}\label{e:next-ips}
\langle v, x_i(k+1)\rangle =
\begin{cases}
(1-\gamma)\langle v, x_i(k)\rangle + \gamma \langle v, x_M(k)\rangle
& \text{if }\langle v, x_i(k)\rangle>0,\\
(1-\gamma)\langle v, x_i(k)\rangle + \gamma \langle v, x_m(k)\rangle
& \text{if }\langle v, x_i(k)\rangle<0.
\end{cases}
\end{equation}
Property \Cref{eq:(d)} at time $k$ and our assumptions that $\langle v,x_m(0)\rangle < 0 < \langle v,x_M(0)\rangle$ imply that $\langle v,x_m(k)\rangle < 0 < \langle v,x_M(k)\rangle$, so we also have $\mathcal C_M(X(k),A)=\{M\}$ and $\mathcal C_m(X(k),A)=\{m\}$. Then, since $\gamma \in (0,1]$, we conclude from \cref{e:next-ips} that properties \Cref{eq:(c)} and \Cref{eq:(d)} hold at time $k+1$. We also conclude that
\begin{align*}
\langle v,x_i(k+1)\rangle-\langle v,x_M(k)\rangle
&=(1-\gamma)\bigl(\langle v,x_i(k)\rangle-\langle v,x_M(k)\rangle\bigr) <0&\text{if }\langle v,x_i(0)\rangle>0,
\\
\langle v,x_i(k+1)\rangle-\langle v,x_m(k)\rangle
&=(1-\gamma)\bigl(\langle v,x_i(k)\rangle-\langle v,x_m(k)\rangle\bigr)>0
&\text{if }\langle v,x_i(0)\rangle<0.
\end{align*}
These inequalities imply that $\langle v,x_m(k+1)\rangle < \langle v,x_i(k+1)\rangle <\langle v,x_M(k+1)\rangle$, which in turn imply that properties \Cref{eq:(a)} and \Cref{eq:(b)} hold at time $k+1$. This concludes the induction process.

Finally, let $n_1 = \bigl|\{i:\langle v,x_i(0)\rangle<0\}\bigr|$, $n_2 = \bigl|\{i:\langle v,x_i(0)\rangle>0\}\bigr|$. For every $i$ with $\langle v,x_i(0)\rangle>0$ we have
$x_i(k+1)-x_M(0)=(1-\gamma)\bigl(x_i(k)-x_M(0)\bigr)$,
and therefore
\[
x_i(k)=x_M(0)+(1-\gamma)^k\bigl(x_i(0)-x_M(0)\bigr)\to x_M(0).
\]
Similar arguments show that $x_i(k)\to x_m(0)$ for every $i$ with $\langle v,x_i(0)\rangle<0$. 
Consequently the limiting configuration consists of $n_1$ copies of $x_m(0)$ and $n_2$ copies of $x_M(0)$ with $n_1+n_2=n$. This is easily seen to be an equilibrium for the dynamical map \cref{eq:hardmaxDynamics}.
\end{proof}

\end{document}
